\documentclass[11pt, a4paper]{article}

\usepackage[utf8]{inputenc}
\usepackage[T1]{fontenc}
\usepackage[english]{babel}
\usepackage{amsmath, amsthm, amssymb}
\usepackage{mathtools}
\usepackage{geometry}
\usepackage{hyperref}

\geometry{margin=2.5cm}

\newtheorem{theorem}{Theorem}[section]
\newtheorem{proposition}[theorem]{Proposition}
\newtheorem{lemma}[theorem]{Lemma}
\newtheorem{corollary}[theorem]{Corollary}
\newtheorem{definition}[theorem]{Definition}
\newtheorem{remark}[theorem]{Remark}
\newtheorem{assumption}{Assumption}[section]

\newcommand{\R}{\mathbb{R}}
\newcommand{\C}{\mathbb{C}}
\newcommand{\E}{\mathbb{E}}
\newcommand{\norm}[1]{\left\|#1\right\|}
\newcommand{\abs}[1]{\left|#1\right|}
\newcommand{\inner}[2]{\left\langle #1,\, #2 \right\rangle}

\title{%
  Bias--Variance Decomposition for\\
  Regularised Laplace Inversion Surrogates
}
\author{}
\date{}

\begin{document}

\maketitle

\begin{abstract}
We derive rigorously the bias--variance decomposition of the total reconstruction
error for a Laplace-domain neural surrogate model combined with a regularised
inversion step. The main result is a bound of the form
\[
  \mathcal{E}_{\mathrm{tot}}
  \;\leq\;
  \frac{\E_\theta\!\left[\norm{(\alpha_t \mathbf{D}_t^*\mathbf{D}_t + \lambda\mathbf{I})\mathbf{v}}\right]}
       {\sigma_{\min}(\mathbf{A})}
  \;+\;
  \frac{\norm{\mathbf{F}}}{\sigma_{\min}(\mathbf{A})}
  \cdot
  \E_\theta\!\left[\norm{\hat{\mathbf{u}} - \hat{\mathbf{u}}^{\mathrm{pred}}}\right],
\]
where $\sigma_{\min}(\mathbf{A})$ denotes the smallest singular value (or equivalently, the smallest eigenvalue) of the normal equation matrix $\mathbf{A}$. The first term is a \emph{regularisation bias} (representing the intrinsic transformation and regularisation error) and the second term is
the \emph{variance} (the surrogate error propagated through the inversion operator).
Both terms are rigorous consequences of the linear structure of the problem.
\end{abstract}

\tableofcontents
\newpage

%==========================================================================
\section{Discrete Setting}
%==========================================================================

\subsection{Notation}

Let $\Omega \subset \R^d$ be the spatial domain, discretised on $N_x$ nodes
$\{x_j\}_{j=1}^{N_x}$.
Time is discretised on a uniform grid $t_n = n\Delta t$, $n = 0,\ldots,N_t-1$,
with time step $\Delta t > 0$ and final time $T = (N_t-1)\Delta t$.

For each simulation parameter $\theta \in \Theta$, the spatio-temporal field
is represented as a matrix
\[
  \mathbf{V}(\theta) \;\in\; \R^{N_x \times N_t},
  \qquad
  V_{j,n}(\theta) = U(x_j, t_n;\theta).
\]
Because the regularised inversion decouples over spatial nodes
(when only temporal regularisation is used), we work with a single
spatial node and drop the index $j$ throughout.
The field at one node is a vector $\mathbf{v}(\theta) \in \R^{N_t}$.

\subsection{Laplace transform matrix}

\begin{assumption}[Coercivity]
  \label{ass:coercivity}
  The underlying spatial operator is strongly elliptic or coercive, meaning there exists a constant $\lambda_1 > 0$ such that the system's eigenvalues have real parts bounded above by $-\lambda_1$. This ensures the field decays exponentially and its Laplace transform is well-defined for $\mathrm{Re}(s) > -\lambda_1$.
\end{assumption}

Fix $N$ complex frequencies $s_1,\ldots,s_N \in \C$ with
$\mathrm{Re}(s_k) > -\lambda_1$ for all $k$, where $\lambda_1 > 0$ is the
coercivity constant of the spatial operator (see Assumption~\ref{ass:coercivity}).

\begin{definition}[Laplace matrix]
  The \emph{Laplace matrix} $\mathbf{F} \in \C^{N \times N_t}$ is defined by
  \begin{equation}
    \label{eq:F_def}
    F_{k,n} \;=\; \Delta t \cdot w_n \cdot e^{-s_k t_n},
    \qquad k = 1,\ldots,N,\quad n = 0,\ldots,N_t-1,
  \end{equation}
  where $w_n$ are quadrature weights
  (trapezoidal: $w_0 = w_{N_t-1} = \tfrac{1}{2}$, $w_n = 1$ otherwise).
\end{definition}

The discrete Laplace transform of $\mathbf{v}$ is then
\begin{equation}
  \label{eq:forward}
  \hat{\mathbf{u}} \;=\; \mathbf{F}\mathbf{v} \;\in\; \C^N.
\end{equation}
Each component $\hat{u}_k = (\mathbf{F}\mathbf{v})_k$ approximates
$\hat{U}(s_k;\theta) = \int_0^T U(t;\theta) e^{-s_k t}\,\mathrm{d}t$.

\subsection{Regularised inversion}

The surrogate predicts $\hat{\mathbf{u}}^{\mathrm{pred}} \in \C^N$.
Given these predictions, we reconstruct $\mathbf{v}^*$ by solving the
regularised least-squares problem
\begin{equation}
  \label{eq:inv_problem}
  \mathbf{v}^*(\theta)
  \;=\;
  \arg\min_{\mathbf{v} \in \R^{N_t}}
  \;\norm{\mathbf{F}\mathbf{v} - \hat{\mathbf{u}}^{\mathrm{pred}}(\theta)}_{\C^N}^2
  + \alpha_t \norm{\mathbf{D}_t \mathbf{v}}^2
  + \lambda \norm{\mathbf{v}}^2,
\end{equation}
where:
\begin{itemize}
  \item $\alpha_t \geq 0$ is a temporal smoothing parameter,
  \item $\lambda > 0$ is a Tikhonov regularisation parameter,
  \item $\mathbf{D}_t \in \R^{(N_t-1)\times N_t}$ is the first-difference
        operator:
        \[
          (\mathbf{D}_t \mathbf{v})_n = v_{n+1} - v_n,
          \quad n = 0,\ldots,N_t-2.
        \]
\end{itemize}

\begin{remark}
  The constraint $\mathbf{v} \in \R^{N_t}$ (real-valued reconstruction) is
  handled by splitting $\mathbf{F} = \mathrm{Re}(\mathbf{F}) + i\,\mathrm{Im}(\mathbf{F})$
  and forming the real augmented system
  \[
    \begin{pmatrix} \mathrm{Re}(\mathbf{F}) \\ \mathrm{Im}(\mathbf{F}) \end{pmatrix}
    \mathbf{v}
    \;\approx\;
    \begin{pmatrix} \mathrm{Re}(\hat{\mathbf{u}}^{\mathrm{pred}}) \\
                    \mathrm{Im}(\hat{\mathbf{u}}^{\mathrm{pred}}) \end{pmatrix}.
  \]
  For clarity of exposition we write the analysis in complex form; the
  real-valued case follows identically by replacing $\C^N$ with $\R^{2N}$.
\end{remark}

%==========================================================================
\section{Normal Equations and Closed-Form Solution}
%==========================================================================

\subsection{The operator $\mathbf{A}$}

The optimality condition (normal equations) for \eqref{eq:inv_problem} is:
\begin{align}
  0 &= \frac{\partial}{\partial \mathbf{v}}
  \left[
    \norm{\mathbf{F}\mathbf{v} - \hat{\mathbf{u}}^{\mathrm{pred}}}^2
    + \alpha_t\norm{\mathbf{D}_t\mathbf{v}}^2
    + \lambda\norm{\mathbf{v}}^2
  \right] \notag \\
  &= 2\,\mathbf{F}^*(\mathbf{F}\mathbf{v} - \hat{\mathbf{u}}^{\mathrm{pred}})
     + 2\alpha_t \mathbf{D}_t^*\mathbf{D}_t\mathbf{v}
     + 2\lambda\mathbf{v}. \notag
\end{align}
Rearranging:
\begin{equation}
  \label{eq:normal}
  \underbrace{
    \left(\mathbf{F}^*\mathbf{F}
    + \alpha_t\mathbf{D}_t^*\mathbf{D}_t
    + \lambda\mathbf{I}\right)
  }_{=:\,\mathbf{A}}
  \mathbf{v}^*
  \;=\;
  \mathbf{F}^* \hat{\mathbf{u}}^{\mathrm{pred}}.
\end{equation}

\begin{proposition}[Invertibility of $\mathbf{A}$]
  \label{prop:invertible}
  If $\lambda > 0$, then $\mathbf{A}$ is symmetric positive definite,
  hence invertible.
\end{proposition}

\begin{proof}
  Each of $\mathbf{F}^*\mathbf{F}$, $\mathbf{D}_t^*\mathbf{D}_t$, and
  $\lambda\mathbf{I}$ is positive semi-definite.
  Since $\lambda > 0$, for any $\mathbf{v} \neq \mathbf{0}$:
  \[
    \mathbf{v}^\top \mathbf{A} \mathbf{v}
    \;\geq\; \lambda\norm{\mathbf{v}}^2 \;>\; 0.
  \]
  Hence $\mathbf{A} \succ 0$ and its smallest singular value (which coincides with its smallest eigenvalue) satisfies $\sigma_{\min}(\mathbf{A}) \geq \lambda > 0$.
\end{proof}

\begin{remark}
  When $\lambda = 0$ and $\alpha_t > 0$, invertibility still holds provided
  $s_k \neq 0$ for all $k$. Indeed, $\ker(\mathbf{A}) \subseteq
  \ker(\mathbf{F}) \cap \ker(\mathbf{D}_t)$; the latter forces $\mathbf{v} = c\mathbf{1}$
  (constant), and $\mathbf{F}(c\mathbf{1}) = 0$ would require
  $\sum_n w_n e^{-s_k t_n} = 0$ for all $k$, which fails when $s_k \neq 0$.
  We keep $\lambda > 0$ for simplicity.
\end{remark}

The unique solution of \eqref{eq:normal} is therefore:
\begin{equation}
  \label{eq:solution}
  \boxed{
    \mathbf{v}^*(\theta)
    \;=\;
    \mathbf{A}^{-1}\mathbf{F}^* \hat{\mathbf{u}}^{\mathrm{pred}}(\theta).
  }
\end{equation}

%==========================================================================
\section{Bias--Variance Decomposition}
%==========================================================================

\subsection{Splitting the error}

We decompose the total error by adding and subtracting the intermediate
quantity $\mathbf{A}^{-1}\mathbf{F}^*\hat{\mathbf{u}}(\theta)$,
where $\hat{\mathbf{u}}(\theta) = \mathbf{F}\mathbf{v}(\theta)$ are the
\emph{exact} Laplace coefficients:

\begin{align}
  \mathbf{v}(\theta) - \mathbf{v}^*(\theta)
  &= \mathbf{v}(\theta)
     - \mathbf{A}^{-1}\mathbf{F}^*\hat{\mathbf{u}}^{\mathrm{pred}}(\theta)
     \notag \\
  &= \underbrace{
       \mathbf{v}(\theta) - \mathbf{A}^{-1}\mathbf{F}^*\mathbf{F}\mathbf{v}(\theta)
     }_{\substack{\text{bias term} \\ \text{(transformation error)}}}
   + \underbrace{
       \mathbf{A}^{-1}\mathbf{F}^*
       \bigl(\hat{\mathbf{u}}(\theta) - \hat{\mathbf{u}}^{\mathrm{pred}}(\theta)\bigr)
     }_{\substack{\text{variance term} \\ \text{(surrogate error)}}}.
  \label{eq:decomp}
\end{align}

\subsection{Simplifying the bias term}

Using the definition $\mathbf{A} = \mathbf{F}^*\mathbf{F}
+ \alpha_t\mathbf{D}_t^*\mathbf{D}_t + \lambda\mathbf{I}$:

\begin{align}
  \mathbf{v} - \mathbf{A}^{-1}\mathbf{F}^*\mathbf{F}\mathbf{v}
  &= \mathbf{A}^{-1}\mathbf{A}\mathbf{v}
     - \mathbf{A}^{-1}\mathbf{F}^*\mathbf{F}\mathbf{v}
     \notag \\
  &= \mathbf{A}^{-1}
     \bigl(\mathbf{A} - \mathbf{F}^*\mathbf{F}\bigr)\mathbf{v}
     \notag \\
  &= \mathbf{A}^{-1}
     \bigl(\alpha_t\mathbf{D}_t^*\mathbf{D}_t + \lambda\mathbf{I}\bigr)
     \mathbf{v}.
  \label{eq:bias_simplified}
\end{align}

The bias depends only on the \emph{true} field $\mathbf{v}$ and the
regularisation operators --- it is independent of the surrogate predictions.

\subsection{Bounding each term}

\begin{lemma}[Bound on the bias]
  \label{lem:bias}
  \begin{equation}
    \label{eq:bias_bound}
    \norm{\mathbf{v} - \mathbf{A}^{-1}\mathbf{F}^*\mathbf{F}\mathbf{v}}
    \;\leq\;
    \frac{\norm{(\alpha_t\mathbf{D}_t^*\mathbf{D}_t + \lambda\mathbf{I})\mathbf{v}}}
         {\sigma_{\min}(\mathbf{A})}.
  \end{equation}
\end{lemma}

\begin{proof}
  From \eqref{eq:bias_simplified} and the sub-multiplicativity of the
  operator norm:
  \[
    \norm{\mathbf{A}^{-1}
          (\alpha_t\mathbf{D}_t^*\mathbf{D}_t + \lambda\mathbf{I})\mathbf{v}}
    \;\leq\;
    \norm{\mathbf{A}^{-1}}
    \cdot
    \norm{(\alpha_t\mathbf{D}_t^*\mathbf{D}_t + \lambda\mathbf{I})\mathbf{v}}.
  \]
  Since $\mathbf{A} \succ 0$,
  $\norm{\mathbf{A}^{-1}} = 1/\sigma_{\min}(\mathbf{A})$,
  which gives \eqref{eq:bias_bound}.
\end{proof}

\begin{lemma}[Bound on the variance]
  \label{lem:variance}
  \begin{equation}
    \label{eq:variance_bound}
    \norm{\mathbf{A}^{-1}\mathbf{F}^*
          \bigl(\hat{\mathbf{u}} - \hat{\mathbf{u}}^{\mathrm{pred}}\bigr)}
    \;\leq\;
    \frac{\norm{\mathbf{F}}}
         {\sigma_{\min}(\mathbf{A})}
    \cdot
    \norm{\hat{\mathbf{u}} - \hat{\mathbf{u}}^{\mathrm{pred}}}_{\C^N}.
  \end{equation}
\end{lemma}

\begin{proof}
  Again by sub-multiplicativity:
  \[
    \norm{\mathbf{A}^{-1}\mathbf{F}^*
          (\hat{\mathbf{u}} - \hat{\mathbf{u}}^{\mathrm{pred}})}
    \;\leq\;
    \norm{\mathbf{A}^{-1}}
    \cdot \norm{\mathbf{F}^*}
    \cdot \norm{\hat{\mathbf{u}} - \hat{\mathbf{u}}^{\mathrm{pred}}}.
  \]
  Since $\norm{\mathbf{F}^*} = \norm{\mathbf{F}}$
  (operator norm equals largest singular value, which is the same for
  $\mathbf{F}$ and $\mathbf{F}^*$), we obtain \eqref{eq:variance_bound}.
\end{proof}

\subsection{Main result}

Combining \eqref{eq:decomp}, Lemma~\ref{lem:bias}, and Lemma~\ref{lem:variance}
via the triangle inequality:

\begin{theorem}[Bias--variance bound]
  \label{thm:bv}
  For any $\theta \in \Theta$, the reconstruction error satisfies
  \begin{equation}
    \label{eq:bv_bound}
    \norm{\mathbf{v}(\theta) - \mathbf{v}^*(\theta)}
    \;\leq\;
    \underbrace{
      \frac{\norm{(\alpha_t\mathbf{D}_t^*\mathbf{D}_t + \lambda\mathbf{I})\,\mathbf{v}(\theta)}}
           {\sigma_{\min}(\mathbf{A})}
    }_{\text{regularisation bias}}
    \;+\;
    \underbrace{
      \frac{\norm{\mathbf{F}}}{\sigma_{\min}(\mathbf{A})}
      \cdot
      \norm{\hat{\mathbf{u}}(\theta) - \hat{\mathbf{u}}^{\mathrm{pred}}(\theta)}_{\C^N}
    }_{\text{propagated surrogate error}}.
  \end{equation}
\end{theorem}

Taking expectations over $\theta$ and using Jensen's inequality
($\E[\norm{\cdot}] \leq \sqrt{\E[\norm{\cdot}^2]}$ in the simpler
$L^1$ form):

\begin{corollary}[Expected error bound]
  \label{cor:expected}
  \begin{equation}
    \label{eq:expected_bv}
    \mathcal{E}_{\mathrm{tot}}
    \;\leq\;
    \frac{\E_\theta\!\left[\norm{(\alpha_t\mathbf{D}_t^*\mathbf{D}_t
                                  + \lambda\mathbf{I})\,\mathbf{v}(\theta)}\right]}
         {\sigma_{\min}(\mathbf{A})}
    \;+\;
    \frac{\norm{\mathbf{F}}}{\sigma_{\min}(\mathbf{A})}
    \cdot
    \E_\theta\!\left[\norm{\hat{\mathbf{u}}(\theta)
                           - \hat{\mathbf{u}}^{\mathrm{pred}}(\theta)}_{\C^N}\right].
  \end{equation}
\end{corollary}

%==========================================================================
\section{Interpretation}
%==========================================================================

\subsection{Role of each term}

Inequality \eqref{eq:expected_bv} separates the total error into two
independent contributions.

\paragraph{Regularisation bias.}
The numerator $\norm{(\alpha_t\mathbf{D}_t^*\mathbf{D}_t + \lambda\mathbf{I})\mathbf{v}}$
measures how much the true solution deviates from the prior encoded by the
regularisation:
\begin{itemize}
  \item $\norm{\mathbf{D}_t\mathbf{v}}^2 = \sum_n (v_{n+1}-v_n)^2$ is large when
        $\mathbf{v}$ varies rapidly in time --- e.g.\ during the explosive
        early transient of the CH$_4$ dispersion.
  \item $\lambda\norm{\mathbf{v}}$ penalises large amplitudes.
\end{itemize}
Increasing $\alpha_t$ or $\lambda$ reduces $\sigma_{\min}(\mathbf{A})^{-1}$
(the denominator grows) but simultaneously increases the numerator.
This is the standard \emph{bias--variance trade-off}.

\paragraph{Propagated surrogate error.}
The factor $\norm{\mathbf{F}}/\sigma_{\min}(\mathbf{A})$ is the
\emph{condition number of the inversion} with respect to the surrogate error.
It quantifies how errors in the predicted Laplace coefficients
$\hat{\mathbf{u}}^{\mathrm{pred}}$ are amplified into errors on the
reconstructed field $\mathbf{v}^*$.

Two observations:
\begin{enumerate}
  \item $\norm{\mathbf{F}}$ depends on the choice of frequencies
        $\{s_k\}$ --- it equals the largest singular value of the
        Laplace matrix and is directly controlled by the contour geometry.
  \item $\sigma_{\min}(\mathbf{A}) \geq \sigma_{\min}(\mathbf{F}^*\mathbf{F})
        + \alpha_t\sigma_{\min}(\mathbf{D}_t^*\mathbf{D}_t) + \lambda$,
        so regularisation improves the condition number.
\end{enumerate}

\subsection{What remains to be bounded}

Theorem~\ref{thm:bv} is fully rigorous. The only term that is not yet
characterised is $\E_\theta[\norm{\hat{\mathbf{u}}(\theta) -
\hat{\mathbf{u}}^{\mathrm{pred}}(\theta)}]$ --- the mean surrogate error
on the Laplace coefficients. Bounding this requires additional assumptions
on the surrogate model (autoencoder + MLP), which are stated separately as
working hypotheses and validated empirically.

\begin{remark}
  The bound \eqref{eq:bv_bound} holds for \emph{any} set of frequencies
  $\{s_k\}$ and any regularisation parameters $(\alpha_t, \lambda)$,
  provided $\lambda > 0$. Optimising over these choices is the subject of future work.
\end{remark}


%==========================================================================
\section{Comparison with the Bromwich--FFT Pipeline}
%==========================================================================
 
We now apply the general framework to the pipeline used in the companion
report~\cite{attarian2026}, which combines a truncated Bromwich contour
with explicit FFT inversion. This allows us to (i)~compute the quantities
$\norm{\mathbf{F}}$ and $\sigma_{\min}(\mathbf{A})$ in closed form, and
(ii)~identify precisely what the truncated FFT does \emph{not} capture
within the bias--variance framework.
 
\subsection{Bromwich contour and FFT structure}
 
The Bromwich contour with damping $\gamma \geq 0$ uses $N_t$ equally-spaced
frequencies
\[
  s_k \;=\; \gamma + i\omega_k,
  \qquad
  \omega_k = \frac{2\pi k}{N_t \Delta t},
  \quad k = 0,\ldots,N_t-1.
\]
The Laplace matrix $\mathbf{F} \in \mathbb{C}^{N_t \times N_t}$ factorises as
\begin{equation}
  \label{eq:F_bromwich}
  \mathbf{F} \;=\; \mathbf{D}_{\mathrm{FFT}} \cdot \mathbf{W},
\end{equation}
where $\mathbf{D}_{\mathrm{FFT}}$ is the $N_t\times N_t$ DFT matrix with
$(D_{\mathrm{FFT}})_{k,n} = e^{-i\omega_k t_n}$, and
$\mathbf{W} = \mathrm{diag}(\Delta t\, w_n\, e^{-\gamma t_n})$ encodes the
quadrature weights and the exponential damping.
 
Since $\mathbf{D}_{\mathrm{FFT}}^*\mathbf{D}_{\mathrm{FFT}} = N_t\mathbf{I}$
(Parseval), we obtain
\begin{equation}
  \label{eq:FstarF_bromwich}
  \mathbf{F}^*\mathbf{F}
  \;=\; \mathbf{W}^*\mathbf{D}_{\mathrm{FFT}}^*\mathbf{D}_{\mathrm{FFT}}\mathbf{W}
  \;=\; N_t\,\mathbf{W}^*\mathbf{W}
  \;=\; N_t\,\Delta t^2\,\mathrm{diag}\!\left(w_n^2\, e^{-2\gamma t_n}\right).
\end{equation}
This is \emph{diagonal} --- a key property of the Bromwich contour.
It follows immediately that:
 
\begin{proposition}[Spectral quantities on the full Bromwich contour]
  \label{prop:bromwich_full}
  On the full Bromwich contour with $N = N_t$ frequencies and damping
  $\gamma \in \R$, $\gamma > -\lambda_1$:
  \begin{align}
    \norm{\mathbf{F}}^2
    &= \sigma_{\max}(\mathbf{F}^*\mathbf{F})
     = N_t\,\Delta t^2\,\max_n\!\left(w_n^2\,e^{-2\gamma t_n}\right)
     = N_t\,\Delta t^2\cdot
       \begin{cases}
         1 & \gamma \geq 0, \\
         e^{-2\gamma T} & \gamma < 0,
       \end{cases}
    \label{eq:F_norm_bromwich} \\[6pt]
    \sigma_{\min}(\mathbf{F}^*\mathbf{F})
    &= N_t\,\Delta t^2\,\min_n\!\left(w_n^2\,e^{-2\gamma t_n}\right)
     = N_t\,\Delta t^2\cdot
       \begin{cases}
         e^{-2\gamma T} & \gamma \geq 0, \\
         1 & \gamma < 0.
       \end{cases}
    \label{eq:sigmin_bromwich}
  \end{align}
  In both cases the conditioning is
  \begin{equation}
    \label{eq:cond_bromwich}
    \kappa(\mathbf{F}^*\mathbf{F}) \;=\; e^{2|\gamma| T}.
  \end{equation}
\end{proposition}
 
\begin{remark}[Role of $\gamma$ on the full contour]
  \label{rem:gamma_full}
  Equation~\eqref{eq:cond_bromwich} shows that $\kappa$ is minimised at
  $\gamma = 0$, where $\kappa = 1$ (perfect conditioning). Any deviation
  $\gamma \neq 0$ degrades conditioning exponentially in $T$.
 
  The two signs of $\gamma$ have \emph{asymmetric} effects on the
  variance term $\norm{\mathbf{F}}/\sigma_{\min}(\mathbf{A})$:
  \begin{itemize}
    \item $\gamma > 0$: $\norm{\mathbf{F}} = \sqrt{N_t}\,\Delta t$ is
          unchanged (maximum at $t_0 = 0$), but $\sigma_{\min}$ decreases
          as $e^{-2\gamma T}$, so the ratio grows as $e^{\gamma T}$.
    \item $\gamma < 0$: $\norm{\mathbf{F}}$ grows as $e^{|\gamma|T}$
          (maximum at $t_{N_t-1} = T$), while $\sigma_{\min}$ is unchanged,
          so the ratio again grows as $e^{|\gamma|T}$.
  \end{itemize}
  Hence $\gamma = 0$ also minimises the variance amplification
  $\norm{\mathbf{F}}/\sigma_{\min}$ on the full contour.
  However, as discussed in Section~\ref{sec:trunc} below, the truncated
  contour introduces an additional spectral decay effect of $\gamma$ that
  modifies this conclusion.
\end{remark}
 
\subsection{Truncation to $k_{\max}$ frequencies}
\label{sec:trunc}
 
In practice, the surrogate only predicts the first $k_{\max} \ll N_t$
Laplace modes. Let $\mathbf{F}_{k_{\max}} \in \mathbb{C}^{k_{\max}\times N_t}$
denote the truncated Laplace matrix (first $k_{\max}$ rows of $\mathbf{F}$).
The Gram matrix $\mathbf{F}_{k_{\max}}^*\mathbf{F}_{k_{\max}}$ is no longer
diagonal; it is a weighted Toeplitz matrix:
\begin{equation}
  \label{eq:gram_truncated}
  \left(\mathbf{F}_{k_{\max}}^*\mathbf{F}_{k_{\max}}\right)_{n,n'}
  \;=\; N_t\,\Delta t^2\, e^{-\gamma(t_n+t_{n'})}\,w_n w_{n'}\,
        D_{k_{\max}}(n - n'),
\end{equation}
where $D_{k_{\max}}(m) = \sin(\pi k_{\max} m/N_t)\,/\,\sin(\pi m/N_t)$
is the Dirichlet kernel.
 
\begin{proposition}[Spectral quantities after truncation]
  \label{prop:bromwich_truncated}
  For the truncated Bromwich contour with $k_{\max}$ frequencies and
  damping $\gamma$:
  \begin{align}
    \norm{\mathbf{F}_{k_{\max}}}^2
    &\;\approx\; k_{\max}\,\Delta t^2\,e^{-2\gamma t_{\min}}
     \;=\; k_{\max}\,\Delta t^2 \cdot \begin{cases} 1 & \gamma \geq 0, \\
     e^{-2\gamma T} & \gamma < 0, \end{cases}
    \label{eq:F_norm_trunc} \\[6pt]
    \sigma_{\min}(\mathbf{F}_{k_{\max}}^*\mathbf{F}_{k_{\max}})
    &\;\approx\; k_{\max}\,\Delta t^2\,e^{-2\gamma T} \cdot
      \begin{cases} 1 & \gamma \geq 0, \\ e^{2\gamma T} & \gamma < 0,
      \end{cases}
    \label{eq:sigmin_trunc}
  \end{align}
  so that the conditioning is $\kappa \approx e^{2|\gamma|T}$,
  as on the full contour.
\end{proposition}
 
\begin{remark}[Trade-off induced by $\gamma$ on the truncated contour]
  \label{rem:gamma_trunc}
  The truncation introduces an effect of $\gamma$ that is absent on the
  full contour: it controls the \emph{spectral decay} of the Laplace
  coefficients. The energy of the $k$-th Laplace mode decays as
  $\norm{\hat{u}_k}^2 \sim \norm{U_0}^2\,e^{-2\gamma T}\,\omega_k^{-2}$,
  so:
  \begin{itemize}
    \item $\gamma > 0$ accelerates the decay of $\norm{\hat{u}_k}$ with
          $k$, reducing the spectral truncation error \eqref{eq:trunc_error}
          (fewer modes needed to represent the signal), at the cost of
          worsening conditioning $\kappa = e^{2\gamma T}$.
    \item $\gamma < 0$ slows spectral decay (more energy in high
          frequencies), requiring larger $k_{\max}$ to achieve the same
          truncation error, but improves conditioning.
    \item $\gamma = 0$ provides no spectral acceleration but achieves
          $\kappa = 1$, making the inversion perfectly conditioned.
  \end{itemize}
  The optimal $\gamma$ therefore depends on the relative magnitude of
  the spectral truncation error and the propagated surrogate error,
  and is problem-dependent. This is consistent with the empirical finding
  in the companion report that $\gamma = 0$ is near-optimal for the
  CH$_4$ dataset, where $k_{\max} = 20$ already captures the dominant
  spectral energy.
\end{remark}
 
\begin{remark}[Effect of $k_{\max}$ at fixed $\gamma$]
  Both $\norm{\mathbf{F}_{k_{\max}}}$ and $\sigma_{\min}$ scale as
  $k_{\max}$ at fixed $\gamma$, so their ratio
  \[
    \frac{\norm{\mathbf{F}_{k_{\max}}}}{\sigma_{\min}(\mathbf{A})}
    \;\approx\;
    \frac{1}{\sqrt{k_{\max}}\,\Delta t\,e^{-\gamma T}}
  \]
  \emph{decreases} with $k_{\max}$ --- the variance amplification improves.
  However, the surrogate error $\E_\theta[\norm{\hat{u} -
  \hat{u}^{\mathrm{pred}}}]$ grows with $k_{\max}$, creating a
  trade-off that depends on the surrogate model capacity.
\end{remark}
 
\subsection{The FFT inversion is not the regularised inverse}
 
A critical observation is that the pipeline in the companion report does
\emph{not} solve the regularised least-squares problem
\eqref{eq:inv_problem}. Instead, it applies the explicit FFT inversion
formula:
\begin{equation}
  \label{eq:fft_inv}
  v^*_n \;=\;
  \frac{e^{\gamma t_n}}{\Delta t}
  \sum_{k=0}^{k_{\max}-1} \hat{u}_k^{\mathrm{pred}}\, e^{i\omega_k t_n}
  \;+\;
  \frac{e^{\gamma t_n}}{\Delta t}
  \sum_{k=k_{\max}}^{N_t-1} \bar{U}_k\, e^{i\omega_k t_n},
\end{equation}
where $\bar{U}_k$ is the per-pixel training mean for $k > k_{\max}$.
 
This operation corresponds to $\alpha_t = \lambda = 0$ in
\eqref{eq:inv_problem} \emph{together with an explicit spectral
truncation}. These are two distinct sources of error that do not appear
in the bias--variance decomposition of Theorem~\ref{thm:bv}:
 
\begin{enumerate}
  \item \textbf{Spectral truncation error.}
    Even with a perfect surrogate ($\hat{u}^{\mathrm{pred}} = \hat{u}$),
    the FFT reconstruction omits the high-frequency content:
    \begin{equation}
      \label{eq:trunc_error}
      \mathbf{v} - \mathbf{v}^*_{\mathrm{FFT}}
      \;=\;
      \frac{e^{\gamma t_n}}{\Delta t}
      \sum_{k=k_{\max}}^{N_t-1}
      \bigl(\hat{u}_k - \bar{U}_k\bigr)\, e^{i\omega_k t_n}.
    \end{equation}
    This is the source of the Gibbs oscillations reported in the companion
    paper, which required a post-hoc UNet corrector to suppress.
 
  \item \textbf{Absence of regularisation bias.}
    Setting $\alpha_t = \lambda = 0$ gives zero regularisation bias in
    \eqref{eq:bv_bound}. However, this comes at the cost of the spectral
    truncation error \eqref{eq:trunc_error}, which is not controlled by
    the framework.
\end{enumerate}
 
\subsection{Summary comparison}
 
Table~\ref{tab:comparison} summarises the two pipelines within the
bias--variance framework.
 
\begin{table}[htbp]
  \centering
  \renewcommand{\arraystretch}{1.5}
  \begin{tabular}{p{4.5cm} p{5.5cm} p{5.5cm}}
    \hline
    & \textbf{FFT truncation} & \textbf{Regularised inverse} \\
    & \textbf{(companion report)} & \textbf{(this work)} \\
    \hline
    Inversion formula
      & Explicit: $e^{\gamma t}/\Delta t \cdot \mathrm{FFT}^{-1}$
      & $\mathbf{A}^{-1}\mathbf{F}^*\hat{\mathbf{u}}^{\mathrm{pred}}$ \\
    $\alpha_t,\,\lambda$
      & $0$
      & $> 0$ \\
    Regularisation bias
      & $0$ (within framework)
      & $\alpha_t\norm{\mathbf{D}_t\mathbf{v}}/\sigma_{\min}(\mathbf{A})$ \\
    Spectral truncation error
      & \textbf{Yes} --- source of Gibbs oscillations
      & \textbf{No} \\
    Requires UNet corrector
      & Yes (481k parameters)
      & No \\
    $\norm{\mathbf{F}}$
      & $\sqrt{k_{\max}}\,\Delta t$
      & depends on contour $\Gamma$ \\
    $\sigma_{\min}(\mathbf{A})$
      & $k_{\max}\,\Delta t^2$
      & $\sigma_{\min}(\mathbf{F}^*\mathbf{F}) + \alpha_t\pi^2/N_t^2 + \lambda$ \\
    \hline
  \end{tabular}
  \caption{Comparison of the FFT truncation pipeline and the regularised
           inverse within the bias--variance framework of
           Theorem~\ref{thm:bv}.}
  \label{tab:comparison}
\end{table}
 
The key insight is that the spectral truncation error \eqref{eq:trunc_error}
falls \emph{outside} the bias--variance decomposition of
Theorem~\ref{thm:bv}: it is neither a regularisation bias nor a propagated
surrogate error. It is an intrinsic limitation of the FFT inversion that
the regularised inverse eliminates by construction, at the cost of
introducing a controlled regularisation bias.

 
\begin{thebibliography}{9}
\bibitem{attarian2026}
K.~Attarian,
\textit{Learning Latent Representations for Neural Surrogate Models of
Partial Differential Equations},
\'Ecole Polytechnique / Institut Polytechnique de Paris, April 2026.
\end{thebibliography}

\end{document}
